\documentclass[12pt]{article}      % Specifies the document class

% assorted useful packages
\usepackage{graphicx}
\usepackage{amsmath}
\usepackage{xspace}
\usepackage{url}
\usepackage{moreverb} % provides the boxedverb environment for the file listing

\newcommand{\tap}{\textsf{TAP}\xspace}
%\newcommand{\taptwo}{\mbox{TAP I\hspace{-2pt}I}\xspace}
\newcommand{\gnome}{\textsf{GNOME}\xspace}
\newcommand{\tab}{\texttt{\textbf{<tab>}}} 
\newcommand{\filename}[1]{\texttt{#1}}
\newcommand{\command}[1]{\texttt{#1}}
%\newcommand{\fileline}[1]{\begin{center}\mbox{\texttt{#1}}\end{center}}
%\newcommand{\fileline}[1]{\begin{flushleft}\mbox{\texttt{#1}}\end{flushleft}}
\newcommand{\fileline}[1]{\begin{flushleft}\texttt{#1}\end{flushleft}}
%\newcommand{\fileline}[1]{\begin{flushleft}\verb~{#1}~\end{flushleft}}
%\newcommand{\backslash}{verb+\+}

\title{Setting up \tap for your own location}    % Declares the document's title.
\author{NOAA HAZMAT
\thanks{Hazardous Materials Response Division,
Office of Response and Restoration,
National Oceanic and Atmospheric Administration,
7600 Sand~Point~Way~NE, Seattle WA 98115, USA.\ \ Email:\texttt{ ORR.TAP@noaa.gov.}}}

\begin{document}             % End of preamble and beginning of text.
\maketitle
\tableofcontents
\section{Introduction}


This document provides a reference about the inputs required by the \tap
Viewing Engine. It seeks to provide all the information required to
produce a \tap location with any trajectory model. Please let us know if
any information in this document is missing or unclear, so that we can
improve it.

\section{Background}

\tap is both an application and a methodology for planning. The methodology involves running an oil spill trajectory model many times, driven by historical data, to form an ensemble of possible oil spill trajectories that we hope represent the possible behavior of future potential oil spills. The results of all those model runs are then examined in the \tap Viewing Engine, which generates maps and statistics that can aid in spill response planning. This same data could be analysis and viewed in other ways, such as a GIS.

\subsection{Trajectory Model}
In order to generate the data for TAP, first you need to set up an oil spill trajectory model for the region of interest. GNOME is one such model, and a good choice for TAP, as we can provide guidance on its use. Virtually any other model could be used, with some work on translating output formats.

The other requirement for TAP is a data set of historical winds and currents with which to drive the trajectory model. The goal is to be able to hindcast enough of the winds and currents in your region of interest to capture the range of behavior of possible spills in the region.

Wind data is widely available in many regions. We feel that at least 10 years or so is required. You do need to be careful that the measured dat reflects your region of interest. This can be on issue if the measurements were taken far away or in regions with strong topographic steering of the wind.

Currents in tidally dominated regions can be modeled by a wide variety of oceanographic models. Currents on the open coast can be far more challenging. Coastal currents can be highly variable and driven by complex forcing, often feeling the effects of forcing far from the region of interest.

In order to hindcast such currents, a hydrodynamic model of the region would need to be applied, with a way to provide boundary conditions for a substantial time period in the past. There may be oceanographers that study the region of interest that might be running such a model, or be able to assist you in getting it set up.

Once you've got your wind data and current model, you can set up GNOME, or your trajectory model of choice, to use those currents and your wind data to model the movement of oil in the area. You can then run GNOME in a batch mode to create a set of trajectories from random times in the past, and then process these to create the data the TAP viewer needs.

You may need to do some translating of file formats so that GNOME can read the output from your current model, but we have used GNOME with many of the major ocean current models, so we should be able to help.
 
\section{Directory structure}

The \tap Viewing Engine finds the files it needs by looking in specific places.  In
the same directory as the \tap application, there must be a
directory (folder) called \filename{TAPDATA}.  This directory has all
the data \tap needs to run.  In this directory are three items:
\begin{itemize}
\item A Map file (BNA format).
\item The \filename{SITE.TXT} file. This file includes a large set of 
data about the site.
\item A set of directories, one for each season represented, that hold 
the binary cube data.
\end{itemize}

In addition, there can optionally be a file in the same place as the
executable called: \filename{TAPCONFIG.TXT}. In this file can be
commands that overwrite some of the defaults in the program. 


\subsection{The map}
Each \tap location must have a base map that is used as the 
background for the display. The map file provides a set of polygons 
surrounding the land, in decimal latitude-longitude coordinates, in 
the BNA format (It can also be in the PICT 
format, but that would only work on a Macintosh). See 
section~\ref{sec:BNA} for a description of this file format.

\section{File formats}
\subsection{BNA format} \label{sec:BNA}
The BNA format is used to describe polygons in lat-long
coordinates.  In \tap, it is used for the base map, and to describe
the receptor sites.  The BNA format is a ASCII comma-delimited text
format that lists polygons one by one.  Each polygon is composed of a
header line, and a set of coordinates.  The header line looks like:
\fileline{"Polygon Name","1",num-points} The first two items are inside
double quotes.  The first item is the name of the polygon.  It is not
used in the map file by \tap.  The second item is a code that
indicates something about the polygon.  \tap uses two codes:
\command{"1"} is a land polygon, and \command{"2"} is a water polygon
that is on top of land (lakes).  The third item is an integer (not
quoted) that is the number of points in the polygon.  Each point is
represented by longitude, latitude, in decimal degrees.  The western
hemisphere is negative longitude, and the southern hemisphere is
negative latitude.  The polygons are all drawn as closed (and filled)
polygons. The first and last points should be the same, but if they are
not, \tap will close the polygon. If you want a polyline, rather than
a filled polygon (bridges, etc.) then you need to create a zero area
polygon by using separate single line segments, or doubling back on the
line so that the polygon has zero area. Figure \ref{fig:bna-example} is
a sample of a small BNA file.

\begin{figure*}[tb!]
\centering
\begin{boxedverbatim}
"Carqueniz Bridge","1",2
-122.22566667,38.0648333
-122.225,38.056333
"Pablo Bridge","1",4
-122.40916667,37.933833
-122.452333,37.9385
-122.4798333,37.944333
-122.452333,37.9385
"1741","1",8
-122.095757,37.453915
-122.097054,37.454048
-122.097580,37.454239
-122.096130,37.456726
-122.095413,37.455772
-122.095009,37.455418
-122.094810,37.455116
-122.095757,37.453915
\end{boxedverbatim}
    \caption{A Sample BNA file. In this example, the second item is 
    drawn as a polyline, because the last points doubles back on the 
    second point.}
    \label{fig:bna-example}
\end{figure*}

\subsection{\filename{SITE.TXT}}\label{sec:site.txt}
The \filename{SITE.TXT} file provides most of the data used to set up 
a \tap location. The 
\filename{SITE.TXT} file included with the \tap demo is fairly 
clearly commented. There should be no blank lines in the file.

The header of the file is whitespace delimited, with all text strings in
double quotes. In the \filename{SITE.TXT} file, any information
following a ``\texttt{//}'' is a comment, and ignored by the
program. The first line in the file is:
\fileline{"SAN FRANCISCO BAY"}
it is the name given to the region. The second line is:
\fileline{"SF.BNA" "BNA"}
The name of the map file, and its type. BNA is recommended, Pict is the 
other option.
\fileline{500 SPILLS }
The number of individual spills modeled (individual scenarios). This 
has to be the same number in every cube.
\fileline{2 SEASONS }
The number of seasons modeled. This can be one or more 
\fileline{"Mar - Aug"  "SFSM" "Summer" \\
          "Sept - Feb"  "SF\_W" "Winter"}
The description and name of the seasons. One line for each season.  The
first field is what gets displayed to the user in the pop-up menu.  The
second field is the prefix of the filenames of the cube files, so that,
in this example, the filenames of the Summer cube files are:
\filename{SFSM0001.bin}, \filename{SFSM0002.bin}, and so on.  The third
field is optional, and is the name of the directory that the cube files
are stored in.  Each season is stored in a separate directory.  If this
field is omitted, the directory name is taken to be the same as the
first field (what the user sees).  Note: make sure that this is a valid
filename.
\fileline{10 TIMES}
The number of output times in each cube.
\fileline{120 "5 days" \\
  96 "4 days" \\
  72 "3 days" \\
  60 "2.5 days" \\
  48 "2 days" \\
  36 "1.5 days"\\ 
  24 "1 day" \\
  12 "12 hours" \\
  6 "6 hours" \\
  1 ""}
The times (from the release) of the output data. The first field is 
in hours. The second field is the string that is presented to the user 
on the pop-up menu. If the second filed is an empty string, the data 
will not be available to the user. These times should be presented in
reverse order (largest first).
\fileline{1 AMOUNTS}
The number of preset spill volumes set for the user.
\fileline{1000 barrels }

The preset volumes set for the user (can be more than one. see the above
line. The allowed units can be any of: \command{barrels, gallons, metric tons, or
cubic meters}. The units are not case sensitive.

\fileline{3 LOCS}
The number of preset Levels of Concern provided in a pop-up menu for the user.
\fileline{10 barrels \\
  5 barrels \\
  1 barrels \\}
The volumes preset for the LOC menu. Possible units are the same as above.
\fileline{185 SITES}
The number of receptor sites.
\fileline{"\#1","1",19 \\
  -122.5306473,    37.8214531 \\
  -122.5305634,    37.8162155  \\
  \ldots }
The polygons of the receptor sites. The polygons are in BNA format 
(see Section~\ref{sec:BNA} There should be one polygon for 
each receptor site. By convention they are labeled as ``\#1'', ``\#2'', 
etc., but you could use any text label you want.

After the receptor site polygons is the following line:
\fileline{233 CUBES}
This line indicates the number of data cubes (one for each spill start 
location).
\fileline{-122.410448,37.863448\\
-122.386002,38.021599 \\
-122.462997,37.833302 \\
\ldots}
The following lines are the locations, in comma separated values in decimal degrees 
of Longitude, Latitude, of the spill locations, one for each cube. It 
is assumed that the spill locations are the same for all seasons.

\subsection{The cube files}
The receptor site hit data is stored in binary files known as the 
cube files. There is one cube file for each spill start location, for 
each represented season. We call them ``cubes'' because the data can 
be thought of as a 3-d array of numbers. The axis are:
\begin{description}
\item[Time] One point for each of the output times specified in the 
\filename{SITE.TXT} file. We usually use ten output times.
\item[Sites] One point for each receptor site, in the order specified 
in the \filename{SITE.TXT} file.
\item[Spills] One point for each realization. Each ``spill'' is a 
single modeled realization of an oil spill from the given spill 
location. The differences between the spills are a result of 
different environmental conditions (winds, tides, etc.). We usually 
use 500 spills, which are each started at a different time, randomly 
sampled from the records we have available. Each of the cubes should 
have the same set of start times, to keep the results consistent.
\end{description}
Each cube is then a 3-d array, with each element in the array 
corresponding to the amount of oil that has passed through a given 
receptor site, by a given time, for a given realization. The amount 
recorded is cumulative, so that oil that hits one site, and then 
re-floats and moves on to a another site is counted at both sites, and the amount 
of oil at a site monotonically increases with time. If the same oil 
hits a site, moves away, and then comes back to the same site, it is 
not counted again.

\subsubsection{Numerical details of the cube files}
In order to minimize storage space, each value in the cube is stored 
in a single (eight bit) byte, as an unsigned integer (number between 0 and 255). 
This number is computed from the fraction of the entire amount of spilled oil 
at that site with the following exponential relationship:
\begin{equation}
x = \operatorname{CEIL}(255f^{\frac{1}{1.5}})
\end{equation}
where x is a the number between 0 and 255 that is stored in the cube 
file, $f$ is the fraction of the total spill that has passed through 
the site by the given time, for the given spill, and the $\operatorname{CEIL}$ 
operator is the ``ceiling'' operator which rounds up to the next 
integer. $\operatorname{CEIL}$ is used in order to give a slightly 
more conservative result than ordinary rounding would. This 
exponential relationship is used in order to give better resolution 
for small amounts of oil than large. For example, the distinction 
between one and two barrels impacting a site is likely to be more 
important than the distinction between 50 and 51 barrels impacting a 
site. In a Lagrangian element model, $f$ is the number of LEs that 
hit the site, divided by the total number of LEs used in the model run.

\paragraph{Order of the data in the file}
While we like to think of the cubes as 3-d arrays, and while many computer 
languages let us store and access them that way, the data must be put 
linearly into the file. The sequence of data in the file is what is 
produced by the following pseudo-code:
\begin{verbatim}
for Time in Times:          
    for Site in Sites:
        for Spill in Spills:
           WriteToFile(Cube[Spill,Site,Time])
\end{verbatim}
Each value is a single eight bit byte, as described above.

We use a program written in the Python programming language
(\url{http://www.python.org}) to generate these cubes. Please contact us 
if you would like a copy to adapt for your own use.


\subsubsection{Cube filenames}
The names of the cube files are a combination of a four letter prefix (defined 
in the \filename{site.txt} file) and the number of the cube, padded 
on the left with zeros to four digits, and the \filename{.BIN} 
extension. This results in an ``8.3'' filename, which should be valid 
on any system. By padding the numbers with zeros, the cube files can 
be alphabetically sorted correctly. For example, if the prefix is
defined as ``\command{PRFX}'' in the \filename{SITE.TXT} file, the
cubes would be named: \filename{PRFX0001.BIN}, \filename{PRFX0002.BIN}, etc.

There is a complete set of cube files, one for each start site, for each 
season modeled. A given numbered cube should correspond to the same 
start site in all seasons.

\section{Batch \gnome}

If you want to use \gnome as the oil spill model to generate \tap cubes,
you'll need to perform the following steps:

\begin{enumerate}
\item Setup a \gnome Save file for your region that includes historical
  data for the time period you are covering. This should always include
  wind data, and often also tides, river flow records, and any other
  scaling factors used in the Save file.
  
\item Set up a \gnome command file (or a set of them, if running on
  multiple machines).
  
\item Run \gnome in batch mode to generate your trajectories.
  
\item Run \filename{BuildCubes.py} (small custom program written in the
  Python Programming language: \url{www.python.org}) to generate your
  cubes. It is available from \url{ORR.TAP@noaa.gov}
  
\item Generate a \tap \filename{SITE.TXT} file for the \tap viewer, and
  put the Cubes generated in the appropriate place in the
  \filename{TAPDATA} folder
\end{enumerate}
Setting up the \gnome save file is documented elsewhere.

NOAA/Hazmat has developed a series of Python scripts that will do much
of this for you. Those scripts are documented in Section \ref{sec:TAPScripts}.



\subsection{Running \gnome in Batch Mode}

The \filename{COMMAND.TXT} file allows generation of Splot File series
to be automated. A Command file is a text file with \gnome commands.
The first line of a Command file must be: 

\fileline{[GNOME Command FILE]}

\begin{description}
\item[Primary method] Launching \gnome with a command file in the \gnome\\directory:

If there is a Command file called \filename{COMMAND.TXT} in the \gnome directory
when \gnome is launched, \gnome will execute the commands in the file
and then quit.

\item[Secondary Method]Opening/Executing command files from within \gnome:

The Windows version of \gnome also supports the execution of a command
file through the normal ``Open'' file.  If you open a file which is
recognized as a \gnome command file, \gnome will execute the commands in
the file.  Note that \gnome will not automatically quit in this case.

The Macintosh version of \gnome also supports the execution of a command
file through the normal ``Open'' file, but in order to use this feature
you will have to set the file's creator to \gnome's signature
\command{COSM} and the file's type to \command{.SAV} .  Not having the
file's type be \command{TEXT} make this method somewhat cumbersome.  An
alternative method is to use drag and drop of a command file.  In order
to get the system to allow you to drop the file onto \gnome, you must
make the creator of the text file be \gnome's signature 'COSM' but you
can leave the file's type as \command{TEXT}.  \gnome will execute the
commands in the file, but will not automatically quit.
\end{description}


\subsubsection{Using an Error Log}

Under typical \gnome usage, any Warning, Error or Note alerts would stop
the model and wait for user to hit the OK button.  To prevent this
behavior, you can specify an Error Log file.  If an error log file is
specified, the text of the alert will be appended to the error log file.
In the case of Warning and Note alerts, the command file continues to
execute.  In the case of an Error alert, execution of the command file
is halted.


\subsubsection{The commands}

In typical usage, the first part of your command file will setup the
data for the scenario(s) you wish to run.  There are two ways to do
this:
\begin{enumerate}
\item Use a command to open a \gnome save file.
\item Use a series of commands of the same type used to setup a location
  file.
\end{enumerate}
In the second part of the command file, you would typically include one
or more lines of the \command{RUNSPILL} message which runs a scenario and outputs
the LE Files.

\paragraph{The \command{RUNSPILL} message}

Here is an example message which includes all the required parameters:

\fileline{MESSAGE runSpill;TO model;  startRelTime 19,07,2000,17,00;
  runDurationInHrs 18;timeStepInMinutes 30; numLEs 1000;startRelPos
  123.7667 W 46.21667 N;outputStepInMinutes 60;outputPath
  :Trajectories$\backslash$Start01$\backslash$Time001.bin;}

The \command{outputPath} parameter specifies the relative path of the
generated Trajectory file. Path elements are separated with
Windows-style backslashes ($\backslash$). If the given path is a file, a
\tap Trajectory File (Sec.~\ref{sec:TrajFile}) will be created. If the
given path is a directory, a set of binary LE files will be created, one
for each time step. That format is documented elsewhere. The \tap
Trajectory File format is much more compact, and the only one we use for
\tap.

Optional parameters for the \command{RUNSPILL} message:\\
\command{endRelTime}, \command{ endRelPos}, \command{ pollutantType},
\command{ totalMass}, \command{ massUnits}, \command{ z.}
\begin{itemize}
\item The default for endRelTime is startRelTime.
\item The default for endRelPos is startRelPos.
\item The default for pollutantType is Non-weathering.
\item The default for totalMass is numLEs.
\item The default for massUnits is Barrels.
\item The default for z is 0.0. Note: z is measured in meters.
\end{itemize}
\command{pollutantType}  KEYWORDS:
\begin{itemize}
\item \command{CONSERVATIVE or "Non-Weathering"}
\item \command{BUNKER or "Fuel Oil \#6"}
\item \command{MEDIUMCRUDE or "Medium Crude"}
\item \command{IFO or "Fuel Oil \#4"}
\item \command{DIESEL}
\item \command{JP4 or "Kerosene / Jet Fuels"}
\item \command{GAS or "Gasoline"}
\end{itemize}

When there are two keywords for the same oil, the first is the one used
in the MOSS files and the second is the one used in the \gnome interface.
You may use either one in the command file.  Note: These key words are
not case sensitive.

\command{massUnits}  KEYWORDS:
\begin{itemize}
\item \command{BARRELS}
\item \command{GALLONS}
\item \command{CUBICMETERS}
\item \command{KILOGRAMS}
\item \command{METRICTONS}
\item \command{SHORTTONS}
\end{itemize}
Additional Messages:
\begin{itemize}
\item MESSAGE close; TO model; 
\item MESSAGE clearWinds; TO model;
\item MESSAGE clearSpills; TO model;
\item MESSAGE reset; TO model; 
\end{itemize}
\subsubsection{An example \filename{COMMAND.TXT} file}

\small{
\begin{verbatim}
[GNOME COMMAND FILE]

-- this would clear all maps etc in case we are not launching GNOME 
MESSAGE close; TO model; 

---------------------------------
-- Set the error log file
---------------------------------
MESSAGE setField; TO model; errorLog :MyFileName.txt;

---------------------------------
-- Either open a save file
---------------------------------
-- MESSAGE open; TO model; PATH yourRelativePathHere;

---------------------------------
-- Or use normal messages from wizard commands in location file
---------------------------------

MESSAGE setfield;TO model;timeStep  0.10;
--
MESSAGE createMap;TO model;TYPE vector; NAME CR Estuary Map; 
PATH :Columbia River:Crest.bna;
--
MESSAGE createMover;TO CR Estuary Map;TYPE Cats;
NAME River Flow.CUR; PATH :Columbia River:River Flow.cur;
--
MESSAGE createMover;TO CR Estuary Map;TYPE Cats; NAME CR Tides.CUR;
PATH :Columbia River:CR Tides.cur
MESSAGE setfield;TO CR Tides.CUR; scaleType constant;
refP 123.7667 W 46.21667 N; scaleValue -1;
timeFile :Columbia River:Shio.txt;

MESSAGE createMover;TO Universal Map;TYPE Random; NAME Diffusion;
MESSAGE setfield;TO Diffusion; coefficient 200000;
uncertaintyFactor 2

---------------------------------
-- specific setup message not typically used in LocationFiles
---------------------------------
MESSAGE clearWinds; TO model;
MESSAGE createMover;TO Universal Map;TYPE Wind; NAME Variable Wind;
PATH :Columbia River:WindData.wnd;
MESSAGE setfield;TO Variable Wind; windage 0.02  0.05;
                                 
---------------------------------
-- specific messages for the TAP run
---------------------------------
MESSAGE runSpill;TO model;  startRelTime 19,07,2000,17,00;
runDurationInHrs 18;timeStepInMinutes 30; numLEs 1000;
startRelPos 123.7667 W 46.21667 N;outputStepInMinutes 60;
outputPath :Trajectories\Start01\Time001.bin;

MESSAGE runSpill;TO model;  pollutantType Kerosene / Jet Fuels;
totalMass 5000; massUnits GALLONS;startRelTime 19,08,2000,17,00;
endRelTime 19,08,2000,21,30; runDurationInHrs 18;
timeStepInMinutes 30;numLEs 1000;startRelPos 123.7667 W 46.21667 N;
endRelPos 123.7667 W 46.24N; outputStepInMinutes 60;
outputPath :Trajectories\Start01\Time001.bin;

---------------------------------
-- if we wanted to quit
--MESSAGE quit; TO model; 

\end{verbatim}
}

\subsection{The \tap Trajectory File Format}\label{sec:TrajFile}
Running \gnome in batch mode generates a set of \tap trajectory files,
which are put in the location specified by the \command{outputPath}
parameter to the \command{runSpill} message in the
\filename{command.txt} file. If you are using NOAA's code for reading
these files, you don't need to know the details, but it could be helpful
for debugging, or if you need to customize the code or write your own.

The trajectory file format is a binary format with a text header. This
allows the user to view the header in any text editor, and examine the
parameters used in the run. The data itself is binary, so that it can be
more compact, and be written and read faster.

\subsubsection{The Header}
The header is in ascii text format, with DOS/Windows style line endings
(\verb+\r\n+), regardless of the system on which \gnome was run. This
allows the files to be transfered between systems, and still be read
with any text editor.

The header has a tagged format, with each section having a tag inside
square brackets (\verb+[]+). Folling the section tag is a set of
key:value pairs separated by a colon. Following the \command{[BINARY
  DATA]} (and a \verb+\r\n+) is the binary data itself. The sections and
data are fairly self explainatory, and best illustrated with an example:

\begin{verbatim}
[FILE INFORMATION]
  File type: TAPRUN
  File version: 1
[RUN PARAMETERS]
  Model start time: January 29, 2000 23:00
  Run duration: 240 hours
  Output time step in seconds: 720
  Number of output steps: 1201
  Number of LEs: 1000
  GNOME version: 1.2.5
  GNOME platform: Windows
  Note:
[BINARY FORMAT]
  Endian: little
  Longitude: LONG
  Latitude: LONG
  Bit flags: CHAR notReleased beached offMaps evaporated notOnSurface
[BINARY DATA]
(binary data here)
\end{verbatim}

Most the key:tag pairs should be clear in meaning from the example. The
\command{[BINARY FORMAT]} section is defined as follows:

\command{Endian:} indicates the endian format of the binary data. For
example, PPC processors uses big endian, and Intel processors use little
endian. If the file is processed on a different type of machine than
\gnome was run on, the data will need to be byteswapped.

The rest of the data describes the structure of the binary data. For
each LE, for each time step, the following data is stored: Latitude,
Longitude, assorted flags. Latitude and Longitude are stored as
integers, representing micro degrees. To convert to floating point
degrees, divide by \command{10e-6}. This format allows us to get plenty
of resolution with minimal storage. Each structure totals 9 bytes:
\begin{itemize}
\item Longitude as a 4 byte C long
\item Latitude as a 4 byte C long
\item Flags as a 1 byte C char
\end{itemize}
The meaning of the flags is specified in order, so that:\\
\verb#flag & 1# is true means that the LE has not been released\\
\verb#flag & 2# is true means the LE is beached\\
etc.\\
Where \verb#&# is the bitwise and operator.

The struct for each LE is written out in sequence, all LEs for each time
step, followed by the next timestep, etc. There are therefore a total of
9 X Number of LEs X Number of Output Steps bytes of binary data.

\section{\tap Scripts for Automating Batch \gnome} \label{sec:TAPScripts}

NOAA/Hazmat has developed a series of scripts, written in the Python
language (\url{www.python.org}), that facilitate setting up \gnome command
files and building the \tap \filename{site.txt} file and \tap cubes.

\subsection{Installation}

\subsubsection{The Python interpreter}

The Python interpreter must be installed to run these scripts. Python is
open source and freely available. On Mac OS-X and many Unix systems, it
comes with the system software. On Windows it must be downloaded and
installed separately. There are a number of sources of information about
Python on the Internet, beginning with the Python home page:
\url{www.python.org}. That page is the best source for documentation and
information about the language. There are also many published books
about Python. While, strictly speaking, you don't need to know anything
about the language to run the scripts, a basic understanding of the
syntax will be helpful, as some scripts do need modification.

One of the strengths of Python is it's extensibility. One of the
extensions available is called numpy, and it is used heavily in
the Hazmat code, so it must also be installed.

\paragraph{Windows}
There are a number of sources of installers for Python for Windows,
starting with the Python home page. You can also install Python and
the Numeric extension separately from:
\begin{itemize}
\item \url{http://www.python.org/download/}
\item \url{http://www.scipy.org/Download}
\end{itemize}
You should use the latest stable versions (Python2.5.2 and
Numpy 1.0.4 as of this writing).

\paragraph{OS-X}

Apple ships recent versions of OS-X with the Python interpreter
included. You should use at least OS-X version 10.4, as previous
versions had python configured differently, and used an older version. With OS-X version 10.5 (Leopard), you can use the built-in python. ON 10.4, you shoould use the installer and associated libraries found here:

\begin{itemize}
  \item \url{http://pythonmac.org/packages/py25-fat/index.html}
\end{itemize}

\paragraph{Unix / Linux}
If a recent version of Python was not supplied with your system, it, and
the numpy package, can be downloaded from the Python and numpy web sites:
\begin{itemize}
  \item \url{http://www.python.org/download/}
  \item \url{http://www.scipy.org/Download}
\end{itemize}
Get the most recent stable versions.

\subsubsection{Editor}
You will need a programmers text editor to edit the config files and
python code. Windows Notepad or Mac TextEdit will do in a pinch, but a
more full featured editor is a good idea. It is particularly nice of it
has a Python mode. If you are already have a favorite, the chances are
good that there is a Python mode available. If you don't currently have
a favorite programmers editor, SciTe is a good, simple choice for
Windows or *nix:
\begin{itemize}
\item \url{http://www.scintilla.org/SciTE.html}
\end{itemize}
On Mac OS-X, TextWrangler is a good, free choice. Jedit, Peppy, or Komodo are other good
choices. 

\subsection{Installing HazmatPy}
Before Using the \tap scripts, a library of routines must be installed
into the python installation on your system. Most of the code in these
routines is straight Python code, and is just ready to use. However, a
few routines are written in C and must be compiled for your system.

\subsubsection{Windows Binary Version}
As Windows does not provide a standard ``system compiler'', compiling the
extensions yourself can be a bit difficult. If you're lucky, you may be
able to simply install a binary distribution. It will look something
like this:
\fileline{Hazmat-1.0.win32-py2.3.exe}
It is a Windows installer. It needs to have Python already installed to
work. 

If you are running a different version of Python that the binary version
was compiled against (Python2.3 in the above example), you can compile
it yourself from the HazmatPy source distribution, as directed below.

\subsubsection{Compiling HazmatPy}
\paragraph{Windows: MSVC}
Windows does not come with a system compiler, so you must have a
compiler installed in order to compile the extensions. The Binary Python
distributions recommended above are compiled with MS VisualC++ v7.  As far as I can tell, extensions must be compiled with the same version of MSVC, that is, version 7.0. The
command line tools must be properly installed for MSVC.

\paragraph{Windows: MinGW}
MSVC is the easiest option on Windows (if you have the right version).
However, an alternative is to use the free MinGw compiler
(\url{http://www.mingw.org}) instead. The process
for doing this can be found at:\\
\url{http://www.mingw.org/MinGWiki/index.php/Python%20extensions}

Once MinGw is installed and working, and the python lib has been patched
as suggested, you can build the package with:
\fileline{python setup.py build --compiler=mingw32}
and install it with:
\fileline{python setup.py install --skip-build}
A binary installer for windows can be created with:
\fileline{python setup.py bdist\_wininst --skip-build}



\paragraph{Linux/ Unix}
If your Python installation was built for your system, it should be all
set to go with the native compiler, usually gcc.

\paragraph{Mac OS-X}
You will need the latest version of the complete OS-X Developer Tools
installed. I think they are now called XCode tools. They can be
downloaded for free from the Apple Developers Connection web site.

\paragraph{All platforms}
Python comes with a built in system for configuring, compiling, and
installing extensions to the python interpreter, called distutils. It
makes it very easy to compile and install packages on a properly
configure system.
\begin{enumerate}
\item Start up a command line interpreter (``DOS box'' on Windows,
  ``Terminal'' on OS-X).
\item change directory to the \filename{HazmatPy} directory.
\item type \fileline{python setup.py build} There may be a few warnings,
  but it should build successfully.
\item To install, type \command{python setup.py --install}. On OS-X,
  you may need to put a \command{sudo} command before the python
  command, as it is being installed into system directories.
\end{enumerate}


\subsection{The Hazmat \tap Scripts}

Hazmat has developed a series of Python scripts to help run \tap. They
are:
\begin{description}
\item[BuildStartTimes.py] a script that examines a set of time series files
  used by \gnome, and computes a set of valid start times that do not
  overlap with any of the ``holes'' in the time series file.
\item[BuildCommandfiles.py] A script that computes a set of command
  files for Batch \gnome.
\item[BuildCubes.py] A script that builds a set of ``cubes'' for the
  \tap viewer from the trajectories output by \gnome.
\item[BuildSite.txt.py] A script that builds the \filename{site.txt}
  file for the \tap viewer.
\end{description}

All of these scripts load all the necessary config data from the single
config file: \filename{TAP\_Setup.py}. This is a python file, rather
than a regular text file for two reasons:
\begin{enumerate}
\item A parser didn't have to be written for a custom text file format.
\item You can use the full Python syntax and language to specify data
  structures and do computations in the file.
\end{enumerate}
For example, if you have lat, long data in (degrees, minutes) format, rather
than decimal degrees, you can simply put, for example:
\fileline{lat = 38 + 12/60.0}
and the conversion is done for you when the file is imported.

\subsubsection{The TAP\_Setup.py file}
The \filename{TAP\_Setup.py} contains all the information required to
set up and run \gnome and the various scripts required to set up a TAP.
There should be no need to edit any other files. The file is full of
comments (denoted by a \command{\#}), and examples that should serve to
explain the contents. Some familiarity with Python syntax and the
workings of \tap will be helpful.

\subsubsection{The Procedure}
\begin{enumerate}
\item Edit the \filename{TAP\_Setup.py} file to suite your situation. The
  scripts should have been delivered with a test package that should be
  all set to go. The file is fairly well commented. When in doubt, leave
  the defaults as they are set, you can always go back and change
  anything later. Also, it's a good idea to comment out existing
  settings, rather than delete them (to a point), so that you retain the
  examples. A Python-aware editor is quite helpful for this.

\item Run \command{python BuildStartTimes.py}. This should create a set
  of start times files, one for each season.

\item Run \command{python BuildCommandFiles.py}. This will create a set
  of \filename{Command.txt} files for \gnome. If you are running one
  single processor machine, and want to do it as one big job, then you
  only need one. This is set in \filename{TAP\_Setup.py}, with the
  \command{NumTrajMachines} parameter. 
  
\item Start Up \gnome, and open the command file you want to run, with
  the \mbox{file--open} menu item. You may have to change the file
  select dialog to show you ``all files'', to see the command file.
  \gnome must be in diagnostic mode to run command files. If it is not in
  diagnostic mode, go to the \mbox{file--preferences} menu item, select
  the ``mode'' tab, then click ``diagnostic'' for both options, and then
  click ``OK''.
  
  After selecting the command file, \gnome should start running in batch
  mode. It will only display a little bit of what's going on on the
  screen, but the window will be locked, and can not be minimized or
  moved. You should be able to use other programs, however (at least on
  Windows XP). On Windows, you can run more than one instance of \gnome
  at once if you have the most recent version. This can be helpful if
  you have more than one processor, or for testing.

  When \gnome is done, it will have created a set of trajectory files
  in the location specified in \filename{TAP\-Setup.py}.

\item Run \command{python BuildCubes.py}. This will build \tap cubes
  from the trajectories generated by \gnome and the grid definition (or
  polygon receptor sites) defined in \filename{TAP\-Setup.py}. They will
  be stored in the location specified in \filename{TAP\-Setup.py}.

  Note that multiple copies of \command{python BuildCubes.py} can be
  run at once, and the script will test to see what's already been
  built. This can be useful if you are building many cubes and want to
  make use of multiple machines and/or multiple processors on one machine.

\item Run \command{python BuildSite.txt.py}. This script creates a
  \filename{site.txt} file for the \tap Viewer that matches all the
  parameters in \filename{TAP\-Setup.py}.

\item Run\command{python BuildViewer.py}. This will copy the
  \filename{site.txt} file and Cubes to the appropriate place for the
  \tap Viewer to use them. You must also have \filename{TAP.EXE} in that
  directory.

\item Run \filename{TAP.EXE}. It should run with your new dataset.


\end{enumerate}










 












\end{document}

% LocalWords:  boxedverb NE hazmat noaa gov TAPDATA BNA TXT lat num polyline tb
% LocalWords:  Carqueniz SF Aug SFSM PRFX LOCS LOC LEs WriteToFile txt
% LocalWords:  LocalWords
